\chapter{Weryfikacja rozwiązania}
\label{ch:testy}

W~celu weryfikacji poprawności działania zaprojektowanego systemu przeprowadzono serię testów manualnych obejmujących
wszystkie kluczowe komponenty architektury.
Testy miały na celu potwierdzenie, że każdy z~mikroserwisów realizuje przypisane mu funkcje zgodnie z~założeniami
przedstawionymi w~rozdziałach~\ref{ch:analiza} i~\ref{ch:projekt}.

Scenariusze testowe skupiały się na weryfikacji następujących aspektów systemu:
\begin{itemize}
    \item Poprawność połączenia z~zewnętrznymi źródłami danych (WebSocket Binance) oraz odbierania strumieniowych
    zdarzeń rynkowych.
    \item Funkcjonowanie magistrali komunikacyjnej NATS JetStream, w~tym publikowanie i~subskrybowanie zdarzeń przez
    poszczególne komponenty.
    \item Działanie systemu powiadomień Telegram oraz modułu zapisu do bazy MongoDB\@.
    \item Działanie mechanizmu filtrowania powiadomień w~zależności od parametrów zdarzeń.
    \item Poprawność zapytań GraphQL z~różnymi filtrami, zarówno zwracającymi wyniki, jak i~pustymi zbiorami danych.
    \item Proces ETL oraz mechanizm \textit{tiered storage}, obejmujący przenoszenie danych z~warstwy gorącej
    (MongoDB) do warstwy zimnej (MinIO) oraz tworzenie zbiorów danych w~formacie Parquet.
\end{itemize}

Każdy test został przeprowadzony w~środowisku Docker z~wykorzystaniem obrazów zbudowanych na podstawie kodu źródłowego
projektu.
Wszystkie scenariusze zakończyły się wynikiem pozytywnym, potwierdzając tym samym zgodność implementacji z~założeniami
projektowymi.

Poniżej przedstawiono szczegółowe opisy poszczególnych scenariuszy testowych.


\newpage


\section[Test T1: Odbiór danych z Binance WebSocket]{Test T1: Odbiór danych z~Binance WebSocket}
\label{sec:test_binance_websocket}

Test weryfikuje poprawność połączenia serwera strumieniowego z~API Binance oraz odbierania komunikatów WebSocket
zawierających zdarzenia typu \textit{aggTrade}.

\vspace{0.5cm}

\begin{table}[H]
    \centering
    \caption{Scenariusz testowy T1: Odbiór danych z~Binance WebSocket.}
    \label{tab:test_binance_websocket}
    \begin{tabular}{|p{0.25\textwidth}|p{0.65\textwidth}|}
        \hline
        \textbf{Pole}        & \textbf{Treść}                                                                                                                                                                                                                                      \\ \hline
        Nazwa scenariusza    & Odbiór danych rynkowych z~Binance WebSocket                                                                                                                                                                                                         \\ \hline
        Cel testu            & Weryfikacja poprawności nawiązania połączenia WebSocket z~API Binance oraz odbierania strumieniowych zdarzeń transakcyjnych dla wybranych par handlowych                                                                                            \\ \hline
        Parametry testu      & Lista symboli: \textit{BTCUSDT, ETHUSDT}; URL WebSocket: \textit{wss://fstream.binance.com/stream?streams=} \newline \textit{btcusdt@aggTrade/ethusdt@aggTrade}                                                                                     \\ \hline
        Kroki testowe        & 1. Uruchomienie kontenera \textit{stream\_server}. \newline 2. Nawiązanie połączenia z~endpointem Binance WebSocket. \newline 3. Nasłuchiwanie wiadomości przychodzących. \newline 4. Analiza logów aplikacji w~celu potwierdzenia odbioru zdarzeń. \\ \hline
        Oczekiwany rezultat  & Aplikacja nawiązuje połączenie WebSocket, odbiera komunikaty w~formacie JSON zawierające pola: \textit{stream}, \textit{data} (z~zagnieżdżonymi polami \textit{e, E, s, p, q, T, m}), oraz loguje odbiór zdarzeń.                                   \\ \hline
        Rzeczywisty rezultat & Połączenie zostało nawiązane poprawnie. Aplikacja otrzymała komunikaty JSON zgodne z~dokumentacją Binance, w~logach pojawiły się wpisy potwierdzające odbiór zdarzeń typu \textit{aggTrade}.                                                        \\ \hline
    \end{tabular}
\end{table}

Test potwierdził, że serwer strumieniowy poprawnie nawiązuje połączenie z~API Binance oraz odbiera dane rynkowe
w~czasie rzeczywistym.
Otrzymane komunikaty zawierały wszystkie wymagane pola domenowe (symbol, cena, wolumen, znacznik czasowy), co
umożliwiło dalsze przetwarzanie zdarzeń przez system.

\newpage


\section{Test T2: Publikowanie zdarzeń do NATS JetStream}
\label{sec:test_nats_publish}

Test weryfikuje poprawność publikowania odebranych zdarzeń z~Binance do magistrali komunikacyjnej NATS JetStream oraz
ich dostępność w~kolejce dla innych konsumentów.

\vspace{0.5cm}

\begin{table}[H]
    \centering
    \caption{Scenariusz testowy T2: Publikowanie zdarzeń do NATS JetStream.}
    \label{tab:test_nats_publish}
    \begin{tabular}{|p{0.25\textwidth}|p{0.65\textwidth}|}
        \hline
        \textbf{Pole}        & \textbf{Treść}                                                                                                                                                                                                                                                                        \\ \hline
        Nazwa scenariusza    & Publikowanie zdarzeń do kolejki NATS JetStream                                                                                                                                                                                                                                        \\ \hline
        Cel testu            & Weryfikacja poprawności publikowania zdarzeń typu \textit{BinanceTradeEvent} do strumienia NATS oraz ich trwałego zapisu w~kolejce                                                                                                                                                    \\ \hline
        Parametry testu      & Stream: \textit{EVENTS}, Subject: \textit{input.events}, Typ zdarzenia: \textit{BinanceTradeEvent} zawierający pola: \textit{ID, Subscribers, Type, Source, EventData, CreatedAt}                                                                                                     \\ \hline
        Kroki testowe        & 1. Uruchomienie serwera strumieniowego oraz NATS JetStream. \newline 2. Odebranie komunikatu z~Binance WebSocket. \newline 3. Opakowanie danych w~strukturę \textit{Event} i~publikacja do kolejki. \newline 4. Weryfikacja obecności zdarzenia w~streamie za pomocą logów aplikacji. \\ \hline
        Oczekiwany rezultat  & Zdarzenie zostaje poprawnie opublikowane do strumienia \textit{EVENTS} na subject \textit{input.events}. W~logach pojawia się wpis potwierdzający publikację wraz z~identyfikatorem zdarzenia. Zdarzenie jest dostępne dla konsumentów subskrybujących dany strumień.                 \\ \hline
        Rzeczywisty rezultat & Zdarzenie zostało pomyślnie opublikowane do NATS JetStream. Logi aplikacji potwierdziły publikację z~odpowiednim identyfikatorem UUID\@. Zdarzenie znalazło się w~kolejce i~jest dostępne dla konsumentów.                                                                            \\ \hline
    \end{tabular}
\end{table}

Test potwierdził poprawne działanie mechanizmu publikowania zdarzeń do magistrali NATS JetStream.
Serwer strumieniowy przekształca surowe komunikaty z~WebSocket w~strukturę \textit{Event} zawierającą metadane (typ,
źródło, lista subskrybentów), a~następnie publikuje je do kolejki.

\newpage


\section{Test T3: Subskrybowanie zdarzeń przez moduł MongoDB}
\label{sec:test_mongo_subscribe}

Test weryfikuje poprawność subskrybowania zdarzeń z~kolejki NATS JetStream przez moduł MongoDB oraz zapisywania ich do
odpowiednich kolekcji.

\vspace{0.5cm}

\begin{table}[H]
    \centering
    \caption{Scenariusz testowy T3: Subskrybowanie zdarzeń przez moduł MongoDB.}
    \label{tab:test_mongo_subscribe}
    \begin{tabular}{|p{0.25\textwidth}|p{0.65\textwidth}|}
        \hline
        \textbf{Pole}        & \textbf{Treść}                                                                                                                                                                                                                                                                                           \\ \hline
        Nazwa scenariusza    & Odbiór zdarzeń przez konsument MongoDB                                                                                                                                                                                                                                                                   \\ \hline
        Cel testu            & Weryfikacja poprawności subskrybowania zdarzeń z~kolejki NATS oraz zapisywania ich do bazy MongoDB w~odpowiedniej kolekcji                                                                                                                                                                               \\ \hline
        Parametry testu      & Subject: \textit{consumer.mongo}, Durable: \textit{mongo-subscriber-durable}, Typ zdarzenia: \textit{BinanceTradeEvent}, Kolekcja docelowa: \textit{BinanceTradeEvent}                                                                                                                                   \\ \hline
        Kroki testowe        & 1. Uruchomienie kontenerów NATS oraz MongoDB. \newline 2. Publikacja zdarzenia do kolejki NATS na subject \textit{consumer.mongo}. \newline 3. Oczekiwanie na przetworzenie zdarzenia przez konsumenta. \newline 4. Weryfikacja obecności rekordu w~bazie MongoDB w~kolekcji \textit{BinanceTradeEvent}. \\ \hline
        Oczekiwany rezultat  & Konsument MongoDB odbiera zdarzenie z~kolejki NATS, wyodrębnia pola: \textit{event\_id, subscriber, type, source, eventData, createdAt}, tworzy dokument BSON i~zapisuje go do kolekcji. W~logach pojawia się wpis potwierdzający zapis z~identyfikatorem zdarzenia.                                     \\ \hline
        Rzeczywisty rezultat & Zdarzenie zostało odebrane i~przetworzone poprawnie. Dokument z~odpowiednimi polami został zapisany do kolekcji \textit{BinanceTradeEvent}. Logi potwierdziły operację z~identyfikatorem zdarzenia.                                                                                                      \\ \hline
    \end{tabular}
\end{table}

Test potwierdził poprawne działanie konsumenta MongoDB jako subskrybenta zdarzeń z~magistrali NATS\@.
Moduł poprawnie odbiera komunikaty, deserializuje je z~formatu JSON do struktury \textit{Event}, konwertuje na
dokumenty BSON oraz zapisuje do odpowiednich kolekcji na podstawie pola \textit{type}.

\newpage


\section{Test T4: Subskrybowanie zdarzeń przez moduł Telegram Bot}
\label{sec:test_telegram_subscribe}

Test weryfikuje poprawność subskrybowania zdarzeń z~kolejki NATS JetStream przez moduł Telegram Bot oraz przetwarzania
komunikatów w~celu wysyłania powiadomień.

\vspace{0.5cm}

\begin{table}[H]
    \centering
    \caption{Scenariusz testowy T4: Subskrybowanie zdarzeń przez moduł Telegram Bot.}
    \label{tab:test_telegram_subscribe}
    \begin{tabular}{|p{0.25\textwidth}|p{0.65\textwidth}|}
        \hline
        \textbf{Pole}        & \textbf{Treść}                                                                                                                                                                                                                                                                                                          \\ \hline
        Nazwa scenariusza    & Odbiór zdarzeń przez konsument Telegram Bot                                                                                                                                                                                                                                                                             \\ \hline
        Cel testu            & Weryfikacja poprawności subskrybowania zdarzeń z~kolejki NATS przez moduł Telegram Bot oraz przetwarzania komunikatów                                                                                                                                                                                                   \\ \hline
        Parametry testu      & Subject: \textit{consumer.telegrambot}, Durable: \textit{telegram-bot-durable}, Typ zdarzenia: \textit{BinanceTradeEvent}                                                                                                                                                                                               \\ \hline
        Kroki testowe        & 1. Uruchomienie kontenerów NATS oraz \textit{telegram\_bot}. \newline 2. Publikacja zdarzenia do kolejki NATS na subject \textit{consumer.telegrambot}. \newline 3. Oczekiwanie na przetworzenie zdarzenia przez konsument. \newline 4. Analiza logów aplikacji w~celu potwierdzenia odbioru i~przetworzenia zdarzenia. \\ \hline
        Oczekiwany rezultat  & Konsument Telegram Bot odbiera zdarzenie z~kolejki NATS, deserializuje je do struktury \textit{Event}, wyodrębnia dane \textit{BinanceTradeData} i~przetwarza je zgodnie z~logiką biznesową. W~logach pojawia się wpis potwierdzający odbiór zdarzenia.                                                                 \\ \hline
        Rzeczywisty rezultat & Zdarzenie zostało odebrane i~przetworzone poprawnie. Moduł Telegram Bot zdeserializował komunikat, wyodrębnił dane transakcji. Logi potwierdziły operację z~odpowiednim identyfikatorem zdarzenia.                                                                                                                      \\ \hline
    \end{tabular}
\end{table}

Test potwierdził poprawne działanie konsumenta Telegram Bot jako subskrybenta zdarzeń z~magistrali NATS\@.
Moduł poprawnie odbiera komunikaty oraz deserializuje zagnieżdżone dane typu \textit{BinanceTradeData} z~pola
\textit{EventData}, przygotowując je do dalszego przetwarzania przez mechanizm filtrowania.

\newpage


\section{Test T5: Wysyłanie powiadomienia Telegram -- filtr spełniony}
\label{sec:test_telegram_filter_pass}

Test weryfikuje poprawność działania mechanizmu filtrowania zdarzeń oraz wysyłania powiadomień przez Telegram Bot
w~przypadku spełnienia warunku filtra.

\vspace{0.5cm}

\begin{table}[H]
    \centering
    \caption{Scenariusz testowy T5: Wysyłanie powiadomienia Telegram -- filtr spełniony.}
    \label{tab:test_telegram_filter_pass}
    \begin{tabular}{|p{0.25\textwidth}|p{0.65\textwidth}|}
        \hline
        \textbf{Pole}        & \textbf{Treść}                                                                                                                                                                                                                                                                                                                                                       \\ \hline
        Nazwa scenariusza    & Wysyłanie powiadomienia przy spełnieniu warunku filtra                                                                                                                                                                                                                                                                                                               \\ \hline
        Cel testu            & Weryfikacja poprawności działania mechanizmu filtrowania oraz wysyłania powiadomień do Telegram w~przypadku, gdy wartość transakcji przekracza ustalony próg                                                                                                                                                                                                         \\ \hline
        Parametry testu      & Zdarzenie \textit{BinanceTradeEvent} z~polami: \textit{Price = 50000.0}, \textit{Quantity = 50.0}, wartość całkowita transakcji: 2~500~000 USD, próg filtra: 2~000~000 USD                                                                                                                                                                                           \\ \hline
        Kroki testowe        & 1. Publikacja zdarzenia do kolejki NATS na subject \textit{consumer.telegrambot}. \newline 2. Oczekiwanie na przetworzenie zdarzenia przez Telegram Bot. \newline 3. Weryfikacja warunku filtra (cena × ilość > 2~000~000). \newline 4. Sprawdzenie logów potwierdzających wysłanie powiadomienia. \newline 5. Weryfikacja otrzymania wiadomości na czacie Telegram. \\ \hline
        Oczekiwany rezultat  & Moduł Telegram Bot oblicza wartość całkowitą transakcji (50000.0 × 50.0 = 2~500~000), porównuje ją z~progiem 2~000~000 USD, stwierdza spełnienie warunku i~wysyła powiadomienie do Telegram API zawierające symbol, cenę, ilość oraz sformatowaną wartość całkowitą. W~logach pojawia się wpis o~wysłaniu alertu.                                                    \\ \hline
        Rzeczywisty rezultat & Warunek filtra został spełniony. Telegram Bot wysłał powiadomienie z~odpowiednimi danymi transakcji. Wiadomość została dostarczona na czat Telegram z~prawidłowo sformatowaną wartością całkowitą (2~500~000.00). Logi potwierdziły operację.                                                                                                                        \\ \hline
    \end{tabular}
\end{table}

Test potwierdził poprawne działanie mechanizmu filtrowania zdarzeń oraz integracji z~Telegram API\@.
System poprawnie oblicza wartość całkowitą transakcji, porównuje ją z~progiem zdefiniowanym w~funkcji
\textit{ValidateTrade} oraz wysyła powiadomienie w~formacie Markdown z~czytelnie sformatowaną kwotą.

\newpage


\section{Test T6: Brak powiadomienia Telegram -- filtr niespełniony}
\label{sec:test_telegram_filter_fail}

Test weryfikuje poprawność działania mechanizmu filtrowania w~przypadku, gdy zdarzenie nie spełnia warunku wysyłania
powiadomienia.

\vspace{0.5cm}

\begin{table}[H]
    \centering
    \caption{Scenariusz testowy T6: Brak powiadomienia Telegram -- filtr niespełniony.}
    \label{tab:test_telegram_filter_fail}
    \begin{tabular}{|p{0.25\textwidth}|p{0.65\textwidth}|}
        \hline
        \textbf{Pole}        & \textbf{Treść}                                                                                                                                                                                                                                                                                                                                                                                    \\ \hline
        Nazwa scenariusza    & Brak powiadomienia przy niespełnieniu warunku filtra                                                                                                                                                                                                                                                                                                                                              \\ \hline
        Cel testu            & Weryfikacja poprawności działania mechanizmu filtrowania oraz potwierdzenie braku wysłania powiadomienia do Telegram w~przypadku, gdy wartość transakcji nie przekracza ustalonego progu                                                                                                                                                                                                          \\ \hline
        Parametry testu      & Zdarzenie \textit{BinanceTradeEvent} z~polami: \textit{Price = 30000.0}, \textit{Quantity = 10.0}, wartość całkowita transakcji: 300~000 USD, próg filtra: 2~000~000 USD                                                                                                                                                                                                                          \\ \hline
        Kroki testowe        & 1. Publikacja zdarzenia do kolejki NATS na subject \textit{consumer.telegrambot}. \newline 2. Oczekiwanie na przetworzenie zdarzenia przez Telegram Bot. \newline 3. Weryfikacja warunku filtra (cena × ilość > 2~000~000). \newline 4. Sprawdzenie logów potwierdzających odrzucenie zdarzenia bez wysyłania powiadomienia. \newline 5. Potwierdzenie braku nowej wiadomości na czacie Telegram. \\ \hline
        Oczekiwany rezultat  & Moduł Telegram Bot oblicza wartość całkowitą transakcji (30000.0 × 10.0 = 300~000), porównuje ją z~progiem 2~000~000 USD, stwierdza brak spełnienia warunku i~pomija wysłanie powiadomienia. Zdarzenie zostaje potwierdzone (Ack) bez dalszego przetwarzania. Brak wpisu w~logach o~wysłaniu alertu.                                                                                              \\ \hline
        Rzeczywisty rezultat & Warunek filtra nie został spełniony. Telegram Bot przetworzył zdarzenie, obliczył wartość transakcji, stwierdził brak spełnienia progu i~nie wysłał powiadomienia. Zdarzenie zostało potwierdzone poprawnie. Brak nowej wiadomości na czacie Telegram.                                                                                                                                            \\ \hline
    \end{tabular}
\end{table}

Test potwierdził poprawne działanie logiki filtrowania zdarzeń.
System nie generuje nadmiarowych powiadomień dla transakcji o~wartości poniżej ustalonego progu, co zapobiega
zaspamowaniu użytkownika alertami o~małej istotności oraz redukuje liczbę wywołań Telegram API\@.

\newpage


\section[Test T7: Zapytanie GraphQL z filtrami zwracającymi wyniki]{Test T7: Zapytanie GraphQL z~filtrami zwracającymi wyniki}
\label{sec:test_graphql_results}

Test weryfikuje poprawność działania interfejsu GraphQL oraz mechanizmu filtrowania danych zapisanych w~MongoDB
z~wykorzystaniem warunków zawsze spełnionych.

\vspace{0.5cm}

\begin{table}[H]
    \centering
    \caption{Scenariusz testowy T7: Zapytanie GraphQL z~filtrami zwracającymi wyniki.}
    \label{tab:test_graphql_results}
    \begin{tabular}{|p{0.25\textwidth}|p{0.65\textwidth}|}
        \hline
        \textbf{Pole}        & \textbf{Treść}                                                                                                                                                                                                                                                                                                                                                  \\ \hline
        Nazwa scenariusza    & Zapytanie GraphQL z~filtrami zwracającymi rekordy                                                                                                                                                                                                                                                                                                               \\ \hline
        Cel testu            & Weryfikacja poprawności wykonywania zapytań GraphQL z~wykorzystaniem filtrów czasowych oraz potwierdzenie zwracania poprawnych danych z~bazy MongoDB                                                                                                                                                                                                            \\ \hline
        Parametry testu      & Zapytanie GraphQL typu \textit{trades} z~filtrem: \textit{eventTime.after > "2020-01-01"}, limit: 10, sortowanie po polu \textit{eventTime} malejąco                                                                                                                                                                                                            \\ \hline
        Kroki testowe        & 1. Uruchomienie serwera GraphQL oraz MongoDB. \newline 2. Upewnienie się, że baza MongoDB zawiera rekordy. \newline 3. Wykonanie zapytania GraphQL z~filtrem czasowym zawsze spełnionym (data po roku 2020). \newline 4. Analiza odpowiedzi JSON zwróconej przez API\@. \newline 5. Weryfikacja obecności pól: \textit{ID, symbol, price, quantity, eventTime}. \\ \hline
        Oczekiwany rezultat  & Zapytanie GraphQL zostaje przetworzone, filtr \textit{eventTime.after} przekształcony na zapytanie MongoDB, wykonane zapytanie do bazy danych zwraca rekordy spełniające warunek. Odpowiedź JSON zawiera listę transakcji z~wybranymi polami, posortowaną według \textit{eventTime} malejąco.                                                                   \\ \hline
        Rzeczywisty rezultat & Zapytanie zostało wykonane poprawnie. API GraphQL zwróciło listę 10 rekordów z~polami \textit{ID, symbol, price, quantity, eventTime}. Wszystkie zwrócone rekordy spełniały warunek filtra, dane były posortowane malejąco według czasu zdarzenia.                                                                                                              \\ \hline
    \end{tabular}
\end{table}

Test potwierdził poprawne działanie resolvera GraphQL oraz integracji z~warstwą repozytorium MongoDB\@.
System poprawnie mapuje filtry GraphQL na warunki zapytań MongoDB (w~tym filtry czasowe, liczbowe oraz wyrażenia regularne),
wykonuje sortowanie oraz limitowanie wyników zgodnie z~parametrami zapytania.

\newpage


\section[Test T8: Zapytanie GraphQL z filtrami niezwracającymi wyników]{Test T8: Zapytanie GraphQL z~filtrami niezwracającymi wyników}
\label{sec:test_graphql_empty}

Test weryfikuje poprawność obsługi zapytań GraphQL z~filtrami, które nie są spełnione przez żadne rekordy w~bazie
danych.

\vspace{0.5cm}

\begin{table}[H]
    \centering
    \caption{Scenariusz testowy T8: Zapytanie GraphQL z~filtrami niezwracającymi wyników.}
    \label{tab:test_graphql_empty}
    \begin{tabular}{|p{0.25\textwidth}|p{0.65\textwidth}|}
        \hline
        \textbf{Pole}        & \textbf{Treść}                                                                                                                                                                                                                                                                                                                                                                                     \\ \hline
        Nazwa scenariusza    & Zapytanie GraphQL z~filtrem niezwracającym rekordów                                                                                                                                                                                                                                                                                                                                                \\ \hline
        Cel testu            & Weryfikacja poprawności obsługi zapytań GraphQL z~warunkami, które nie są spełnione przez żadne dane w~bazie oraz potwierdzenie zwrócenia pustej listy bez błędów                                                                                                                                                                                                                                  \\ \hline
        Parametry testu      & Zapytanie GraphQL typu \textit{trades} z~filtrem: \textit{eventTime.before < "2000-01-01"}, limit: 10                                                                                                                                                                                                                                                                                              \\ \hline
        Kroki testowe        & 1. Uruchomienie serwera GraphQL oraz MongoDB. \newline 2. Upewnienie się, że baza MongoDB zawiera rekordy z~datami nowszymi niż rok 2000. \newline 3. Wykonanie zapytania GraphQL z~filtrem czasowym nigdy niespełnionym (data przed rokiem 2000). \newline 4. Analiza odpowiedzi JSON zwróconej przez API\@. \newline 5. Weryfikacja, że odpowiedź zawiera pustą listę bez komunikatów o~błędach. \\ \hline
        Oczekiwany rezultat  & Zapytanie GraphQL zostaje przetworzone poprawnie, filtr \textit{eventTime.before} przekształcony na warunek MongoDB\@. Zapytanie do bazy danych nie zwraca żadnych rekordów. Odpowiedź JSON zawiera pustą listę \textit{trades: []} bez błędów, z~kodem statusu HTTP 200.                                                                                                                          \\ \hline
        Rzeczywisty rezultat & Zapytanie zostało wykonane poprawnie. API GraphQL zwróciło pustą listę rekordów \textit{trades: []}. Brak błędów w~odpowiedzi, status HTTP 200. System poprawnie obsłużył sytuację braku wyników spełniających warunek.                                                                                                                                                                            \\ \hline
    \end{tabular}
\end{table}

Test potwierdził poprawną obsługę przypadków brzegowych przez interfejs GraphQL\@.
System nie generuje błędów w~sytuacji, gdy zapytanie z~poprawnymi parametrami nie zwraca żadnych wyników, co jest
zgodne z~konwencjami REST i~GraphQL (pusta lista vs. kod błędu 404).

\newpage


\section[Test T9: Proces ETL i mechanizm tiered storage]{Test T9: Proces ETL i~mechanizm tiered storage}
\label{sec:test_etl_tiered_storage}

Test weryfikuje poprawność działania procesu ETL oraz mechanizmu przenoszenia danych z~warstwy gorącej (MongoDB) do
warstwy zimnej (MinIO) z~jednoczesnym przetwarzaniem danych do formatu Parquet.

\vspace{0.5cm}

\begin{table}[H]
    \centering
    \caption{Scenariusz testowy T9: Proces ETL i~mechanizm tiered storage.}
    \label{tab:test_etl_tiered_storage}
    \begin{tabular}{|p{0.25\textwidth}|p{0.65\textwidth}|}
        \hline
        \textbf{Pole}        & \textbf{Treść}                                                                                                                                                                                                                                                                                                                                                                                                                                                                                                                                             \\ \hline
        Nazwa scenariusza    & Przenoszenie danych z~MongoDB do MinIO oraz generowanie zbiorów analitycznych                                                                                                                                                                                                                                                                                                                                                                                                                                                                              \\ \hline
        Cel testu            & Weryfikacja poprawności procesu ETL obejmującego: pobranie danych z~MongoDB, zapis do warstwy \textit{raw}, wzbogacenie danych w~warstwie \textit{curated}, generowanie agregacji w~warstwie \textit{analytics} oraz usunięcie danych z~MongoDB po pomyślnym transferze                                                                                                                                                                                                                                                                                    \\ \hline
        Parametry testu      & Zbiór rekordów w~MongoDB (kolekcje: \textit{BinanceTradeEvent}, \textit{NotificationEvent}), MinIO bucket: \textit{data-lake}, warstwy: \textit{raw/trades, curated/trades, analytics/minute, analytics/hourly, analytics/daily}                                                                                                                                                                                                                                                                                                                           \\ \hline
        Kroki testowe        & 1. Uruchomienie MongoDB, MinIO oraz modułu ETL\@. \newline 2. Upewnienie się, że MongoDB zawiera rekordy testowe. \newline 3. Wywołanie procesu ETL. \newline 4. Weryfikacja zapisu plików JSONL (gzip) w~warstwie \textit{raw} MinIO\@. \newline 5. Weryfikacja zapisu plików Parquet w~warstwie \textit{curated} MinIO\@. \newline 6. Weryfikacja zapisu plików Parquet w~warstwie \textit{analytics} (trzy poziomy: minutowy, godzinowy, dzienny). \newline 7. Sprawdzenie kolekcji MongoDB -- rekordy powinny zostać usunięte po pomyślnym transferze. \\ \hline
        Oczekiwany rezultat  & Proces ETL pobiera dane z~MongoDB w~partiach, zapisuje surowe dane w~formacie JSONL (gzip) w~folderze \textit{raw/trades} oraz \textit{raw/notifications}. Następnie tworzy wzbogacone rekordy \textit{ProcessedTrade} i~zapisuje je w~formacie Parquet w~folderze \textit{curated}. Moduł analytics agreguje dane do poziomów czasowych i~zapisuje je w~folderze \textit{analytics}. Po pomyślnym zapisie wszystkich warstw, przetworzone rekordy są usuwane z~MongoDB.                                                                                   \\ \hline
        Rzeczywisty rezultat & Proces ETL został wykonany pomyślnie. Pliki JSONL z~surowymi danymi pojawiły się w~warstwie \textit{raw}. Pliki Parquet z~przetworzonymi danymi zostały zapisane w~warstwie \textit{curated}. Pliki Parquet z~agregacjami na trzech poziomach czasowych pojawiły się w~folderach \textit{analytics/minute}, \textit{analytics/hourly}, \textit{analytics/daily}. Kolekcje MongoDB zostały opróżnione.                                                                                                                                                      \\ \hline
    \end{tabular}
\end{table}

Test potwierdził poprawne działanie pełnego łańcucha przetwarzania danych w~modelu \textit{tiered storage}.
Proces ETL realizuje wszystkie założone etapy: ekstrakcję z~warstwy gorącej, transformację danych (wzbogacanie
o~metadane, agregacja), zapis do warstwy zimnej w~formatach zoptymalizowanych pod kątem archiwizacji oraz ładowanie
danych analitycznych w~strukturze partycjonowanej według symbolu i~przedziałów czasowych.
Mechanizm zapewnia odciążenie bazy operacyjnej MongoDB przy zachowaniu dostępności danych historycznych w~magazynie
obiektowym MinIO\@.
